\section{The $2\frac{25}{42}$-Approximation Algorithm}

The approach utilized by the $2\frac{25}{42}$-approximation algorithm is similar
to earlier methods that rely on overlap maximization. The core structural steps
resemble the generic algorithm, with the critical difference lying in the final
step: instead of computing a second cycle cover, the algorithm employs a
sophisticated overlap approximation algorithm to merge the cycle
representatives.

\subsection{Algorithm Outline}
The cycle representatives $t_C$ are chosen using the exact same optimal rotation
conditions as established in the $2\frac{2}{3}$-approximation algorithm. The
steps are defined as follows:

\begin{enumerate}
    \item \textbf{Initial Distance Graph:} Construct the distance graph $G_S$
    for the set $S$.
    \item \textbf{First Cycle Cover:} Find a minimum weight cycle cover
    $\mathcal{C}_S$ on the graph $G_S$.
    \item \textbf{Optimal Representatives:} For each cycle $C \in
    \mathcal{C}_S$, compute and choose the representative string $t_C$ exactly
    as described in Lemma \ref{lemma:representative_choice}.
    \item \textbf{Overlap Approximation:} Let $T$ be the set of all chosen
    strings $t_C$. Construct a superstring of $T$ using a $\delta$-approximation
    maximum overlap algorithm.
\end{enumerate}

\subsection{Approximation Ratio Analysis}
To establish the final approximation ratio, we first need to define the
relationship between the length of the produced superstring and the maximum
overlap of the set $T$.

\begin{lemma} \label{lemma:apx_overlap}
Let $apx(T)$ be the length of the superstring of $T$ produced by a
$\delta$-approximation maximum overlap algorithm. Then:
\begin{equation}
apx(T) \le opt(T) + (1 - \delta)\textit{maxov}(T)
\end{equation}
\end{lemma}

\begin{proof}
Recall the fundamental relationship that $opt(T) = \sum_{t_i \in T} |t_i| -
\textit{maxov}(T)$. Since the overlap achieved by the $\delta$-overlap
approximation algorithm is by definition at least $\delta \cdot
\textit{maxov}(T)$, the length of the resulting string is bounded by:
\begin{equation}
apx(T) \le \sum_{t_i \in T} |t_i| - \delta \cdot \textit{maxov}(T) = opt(T) + (1 - \delta)\textit{maxov}(T)
\end{equation}
\end{proof}

To utilize this lemma, we need a strict upper bound on the possible maximum
overlap $\textit{maxov}(T)$. In older algorithms, the standard bound was
$\textit{maxov}(T) \le 2 CYC(G_S)$. However, due to the specific choice of cycle
representatives $t_C$ in Step 3, this bound can be significantly tightened.

\begin{lemma} \label{lemma:tight_maxov}
The maximum overlap of the set $T$ is bounded by:
\begin{equation}
\textit{maxov}(T) \le \frac{3}{2} CYC(G_S)
\end{equation}
\end{lemma}

\begin{proof}
Consider the shortest superstring for $T$, and assume that it contains the
strings $t_1, t_2, \dots$ in that specific order. Recall that the strings in $T$
are mutually inequivalent due to the minimality of the cycle cover
$\mathcal{C}_S$. Therefore, by applying the Overlap-Rotation Lemma
\ref{lemma:overlap_rotation}, the overlap between any two adjacent strings in
the superstring satisfies:
\begin{equation}
ov(t_i, t_{i+1}) \le \textit{period}(t_i) + \frac{1}{2}\textit{period}(t_{i+1})
\end{equation}
Summing these overlaps over all adjacent strings $t_i \in T$, and recalling that
the period of a representative string equals the weight of its generating cycle,
we obtain:
\begin{equation}
\textit{maxov}(T) = \sum_{i=1}^{|T|-1} ov(t_i, t_{i+1}) \le \frac{3}{2} \sum_{t_i \in T} \textit{period}(t_i) = \frac{3}{2} CYC(G_S)
\end{equation}
\end{proof}

By combining these bounds and utilizing a $\delta$-approximation algorithm for
directed Max-TSP-Path in the final step, the overall approximation ratio is
established. Table~\ref{tab:maxtsp_history} summarizes the history of such
algorithms and the resulting shortest superstring ratio obtained by substituting
$\delta$ into Lemma~\ref{lemma:apx_overlap} and Lemma~\ref{lemma:tight_maxov}.

\begin{table}[htpb]
    \centering
    \caption{History of approximation-ratio improvements for directed
    Max-TSP-Path, and the resulting Shortest Superstring ratio $2 +
    \frac{3}{2}(1-\delta)$ obtained when the algorithm is plugged into Step 4 of
    this section.}
    \label{tab:maxtsp_history}
    \renewcommand{\arraystretch}{1.3}
    \resizebox{\textwidth}{!}{%
    \begin{tabular}{@{}llcc@{}}
        \toprule
        \textbf{Year} & \textbf{Authors} & \textbf{Max-TSP-Path Ratio $\delta$} & \textbf{Resulting SCS Ratio} \\
        \midrule
        \textbf{1994} & \textbf{Kosaraju, Park, Stein \cite{Kosaraju1994}} & $\mathbf{\frac{38}{63} \approx 0.603}$ & $\mathbf{2\frac{25}{42} \approx 2.596}$ \\
        2004 & Bläser \cite{Blaser2004} & $\frac{8}{13} \approx 0.615$ & $2\frac{15}{26} \approx 2.577$ \\
        2005 & Kaplan et al. \cite{kaplan2005approximation} & $\frac{2}{3} \approx 0.667$ & $2\frac{1}{2} = 2.5$ \\
        2012 & Paluch, Elbassioni, van Zuylen \cite{paluch2012simpler} & $\frac{2}{3} \approx 0.667$ & $2\frac{1}{2} = 2.5$ \\
        2020 & Paluch \cite{paluch2020new} & $\frac{7}{10} = 0.7$ & $2\frac{9}{20} = 2.45$ \\
        \bottomrule
    \end{tabular}%
    }
\end{table}

\begin{theorem}
By employing the $\frac{38}{63}$-overlap approximation algorithm proposed by
Kosaraju in Step 4, the algorithm achieves a $2\frac{25}{42} \approx 2.596$
approximation for the shortest superstring problem.
\end{theorem}

\begin{proof}
By substituting the bounds from Lemma \ref{lemma:optT}, Lemma
\ref{lemma:apx_overlap}, and Lemma \ref{lemma:tight_maxov} into the overlap
approximation formula with $\delta = \frac{38}{63}$, the length of the produced
superstring is bounded by:
\begin{equation}
apx(T) \le opt(T) + \left(1 - \frac{38}{63}\right)\textit{maxov}(T)
\end{equation}
\begin{equation}
apx(T) \le 2opt(S) + \left(\frac{25}{63}\right) \frac{3}{2} CYC(G_S)
\end{equation}
Since $CYC(G_S) \le opt(S)$, this simplifies to:
\begin{equation}
apx(T) \le 2opt(S) + \frac{25}{42} opt(S) = 2\frac{25}{42} opt(S)
\end{equation}
This mathematically guarantees the $2.596$ approximation ratio.
\end{proof}

\subsection{Computational Complexity}
Let $n = |S|$ and let $L = \sum_{s \in S} |s|$ denote the total input length.
Steps 1 and 2 build the distance graph $G_S$ and compute its minimum weight
cycle cover, the latter using the Hungarian algorithm in $O(n^3)$ time
\cite{kuhn1955hungarian, edmondskarp1972, tomizawa1971techniques}. As in the $2\frac{2}{3}$-approximation algorithm, Step
3 computes the representatives $t_C$ in $O(L)$ time overall, by Lemma
\ref{lemma:representative_choice}. Step 4 runs a $\delta$-approximation
algorithm for directed Max-TSP-Path on $T$, where $|T| \le n$. Writing
$T_\delta(n)$ for the running time of this subroutine, the algorithm overall
runs in $O(n^3 + T_\delta(n) + L)$ time.

When Kosaraju's $\frac{38}{63}$-approximation algorithm is used for this step
(Theorem \ref{thm:final_approximation}), each of its three versions relies on
the Hungarian algorithm and Edmonds' blossom algorithm, both running in $O(n^3)$
time, so $T_\delta(n) = O(n^3)$. Substituting this back, the algorithm as a
whole runs in $O(n^3 + L)$ time.
